\documentclass[tikz,border=4pt]{standalone}
\usepackage{tikz}
\usepackage{polymer-paper-preamble}
% Reusable Panel F apparatus motif. The caller owns labels, force arrows,
% panel framing, and whole-panel composition.
\newcommand{\PanelFFloatingCantilever}[1]{%
  \begin{scope}[shift={(#1)}]
    % Stable anchors for semantic overlays and downstream inspection.
    \coordinate (panelFFixedBoundary) at (0,0.34);
    \coordinate (panelFFloatingCantilever) at (-0.26,0);
    \coordinate (panelFDrivenElectrode) at (1.34,-1.40);
    \coordinate (panelFSourceReturn) at (1.74,0.39);
    \coordinate (panelFTrappedCharge) at (-0.22,-0.56);

    % Fixed mechanical boundary and shallow root jaw.
    \draw[cGray!62!black, line width=0.72pt, line cap=round, line join=round]
      (-0.26,0.34) -- (0.26,0.34);
    \foreach \hx in {-0.22,-0.08,0.06,0.20} {
      \draw[cGray!48!black, line width=0.30pt]
        (\hx,0.34) -- ({\hx+0.07},0.42);
    }
    \draw[cGray!62!black, line width=0.72pt, line cap=round]
      (0,0.34) -- (0,0.14);
    \fill[cGray!12, rounded corners=0.22mm]
      (-0.16,0) rectangle (0.16,0.14);
    \draw[cGray!62!black, line width=0.72pt, rounded corners=0.22mm]
      (-0.16,0) rectangle (0.16,0.14);
    \draw[cGray!56!black, line width=0.34pt] (-0.12,0) -- (-0.12,-0.06);
    \draw[cGray!56!black, line width=0.34pt] (0.12,0) -- (0.12,-0.06);

    % Electrically floating charge-trapping cantilever.
    \shade[top color=cAmber!18, bottom color=cAmber!46, rounded corners=0.5mm]
      (-0.14,0) .. controls (-0.02,-0.74) and (-0.49,-1.54) ..
      (-1.37,-2.04) -- (-1.22,-2.20) .. controls (-0.28,-1.72) and
      (0.24,-0.78) .. (0.14,0) -- cycle;
    \draw[cAmber!82!black, line width=0.62pt, rounded corners=0.5mm]
      (-0.14,0) .. controls (-0.02,-0.74) and (-0.49,-1.54) ..
      (-1.37,-2.04) -- (-1.22,-2.20) .. controls (-0.28,-1.72) and
      (0.24,-0.78) .. (0.14,0) -- cycle;
    \draw[cAmber!60!black, line width=0.28pt, opacity=0.42]
      (-0.06,-0.16) .. controls (-0.02,-0.76) and (-0.56,-1.54) .. (-1.29,-2.06);
    \draw[cAmber!60!black, line width=0.28pt, opacity=0.35]
      (0.04,-0.12) .. controls (0.14,-0.82) and (-0.39,-1.72) .. (-1.16,-2.16);

    % Driven electrode and gap guides.
    \shade[left color=cGray!30, right color=cGray!16]
      (1.34,-2.38) rectangle (1.58,-0.02);
    \draw[cGray!82!black, line width=0.62pt]
      (1.34,-2.38) rectangle (1.58,-0.02);
    \draw[cGray!45!black, line width=0.25pt] (1.39,-2.34) -- (1.39,-0.06);
    \foreach \hy in {-2.22,-1.98,-1.74,-1.50,-1.26,-1.02,-0.78,-0.54,-0.30} {
      \draw[cGray!48!black, line width=0.20pt]
        (1.58,\hy) -- (1.34,{\hy-0.05});
    }
    \foreach \yy/\xstart in {-1.88/-0.96,-1.50/-0.67,-1.12/-0.44,-0.74/-0.22} {
      \draw[cGray!23, line width=0.30pt, dash pattern=on 1.3pt off 1.2pt]
        (\xstart,\yy) -- (1.34,\yy);
    }
    \draw[cGray!19, line width=0.22pt, dash pattern=on 1.2pt off 1.5pt]
      (1.08,-0.42) -- (1.34,-0.42);

    % Driven source lead and grounded source return. No lead touches sample.
    \fill[cBlue!5, rounded corners=0.55mm] (0.76,0.40) rectangle (1.24,0.78);
    \draw[cBlue!64!black, line width=0.38pt, rounded corners=0.55mm]
      (0.76,0.40) rectangle (1.24,0.78);
    \node[font=\sffamily\bfseries\fontsize{4.4}{5.2}\selectfont,
          text=cBlue!70!black] at (1.00,0.59) {HV};
    \fill[cBlue!70!black] (1.24,0.50) circle (0.025);
    \draw[cBlue!70!black, line width=0.36pt, rounded corners=0.3mm]
      (1.24,0.50) -- (1.46,0.50) -- (1.46,-0.02);
    \fill[cBlue!70!black] (1.24,0.68) circle (0.025);
    \draw[cBlue!58!black, line width=0.34pt, rounded corners=0.3mm]
      (1.24,0.68) -- (1.74,0.68) -- (1.74,0.39);
    \draw[cBlue!58!black, line width=0.34pt] (1.61,0.39) -- (1.87,0.39);
    \draw[cBlue!58!black, line width=0.34pt] (1.65,0.33) -- (1.83,0.33);
    \draw[cBlue!58!black, line width=0.34pt] (1.69,0.27) -- (1.79,0.27);

    % Flat trapped-charge markers.
    \foreach \cx/\cy/\rr in {-0.22/-0.56/0.075,-0.41/-0.98/0.082,-0.74/-1.48/0.088,-1.12/-1.90/0.082} {
      \fill[cRed!68!black] (\cx,\cy) circle (\rr);
      \draw[cRed!88!black, line width=0.25pt] (\cx,\cy) circle (\rr);
    }
  \end{scope}%
}


\tikzset{
  panel letter/.style={
    font=\sffamily\bfseries\fontsize{9.0}{10.5}\selectfont,
    text=black,
    anchor=west
  },
  panel title/.style={
    font=\sffamily\bfseries\fontsize{7.0}{8.4}\selectfont,
    text=cGray,
    anchor=west
  },
  note/.style={
    font=\sffamily\fontsize{5.8}{7.0}\selectfont,
    text=cGray,
    align=center
  },
  small note/.style={
    font=\sffamily\fontsize{5.1}{6.2}\selectfont,
    text=cGray,
    align=center
  },
  evidence axis/.style={
    draw=cGray,
    line width=0.42pt,
    -{Stealth[length=3.0pt,width=2.2pt]}
  },
  causal/.style={
    draw=cGray!55,
    line width=0.36pt,
    -{Stealth[length=3.0pt,width=2.2pt]}
  },
  force arrow/.style={
    draw=cRed,
    line width=0.72pt,
    -{Stealth[length=4.0pt,width=3.0pt]}
  }
}

\begin{document}
\begin{tikzpicture}[line cap=round, line join=round]
  \path[use as bounding box] (-0.35,-0.35) rectangle (18.65,11.25);

  % A — material identity.
  \node[panel letter] at (0.00,10.83) {A};
  \node[panel title] at (0.36,10.83) {Sulfur-rich polymer};
  \draw[cGray!65, line width=0.62pt]
    (0.35,9.78) .. controls (0.85,10.22) and (1.18,9.35) .. (1.72,9.78)
    .. controls (2.17,10.14) and (2.55,9.46) .. (3.13,9.83);
  \foreach \x/\y in {0.48/9.88,1.00/9.74,1.47/9.70,1.94/9.85,2.43/9.68,2.94/9.78} {
    \fill[cLGray!55] (\x,\y) circle (0.105);
    \draw[cGray!72, line width=0.28pt] (\x,\y) circle (0.105);
  }
  \foreach \x/\y in {0.76/9.55,1.27/10.02,1.70/9.48,2.20/10.03,2.72/9.48} {
    \fill[cAmberSphere] (\x,\y) circle (0.135);
    \draw[cAmber, line width=0.34pt] (\x,\y) circle (0.135);
  }
  \node[small note, anchor=west] at (0.34,9.18)
    {polymer backbone with sulfur-rich motifs};

  % B — qualitative composition series.
  \node[panel letter] at (4.18,10.83) {B};
  \node[panel title] at (4.54,10.83) {Sulfur-content series};
  \foreach \x/\label in {4.52/lower S,6.08/intermediate,7.64/higher S} {
    \path[fill=cLGray!20, draw=cGray!70, line width=0.38pt, rounded corners=1.2pt]
      (\x,9.36) rectangle ++(1.15,0.72);
    \node[small note] at ({\x+0.575},9.12) {\label};
  }
  \foreach \x/\y in {4.72/9.55,5.31/9.83,5.46/9.52} {
    \fill[cAmberSphere] (\x,\y) circle (0.075);
  }
  \foreach \x/\y in {6.22/9.55,6.48/9.84,6.77/9.52,7.02/9.78,7.09/9.46} {
    \fill[cAmberSphere] (\x,\y) circle (0.075);
  }
  \foreach \x/\y in {7.78/9.49,8.02/9.83,8.25/9.58,8.52/9.90,8.63/9.49,8.35/9.75,7.91/9.68} {
    \fill[cAmberSphere] (\x,\y) circle (0.075);
  }
  \draw[causal] (4.66,8.78) -- (8.58,8.78);
  \node[small note, fill=white, inner sep=1.2pt] at (6.62,8.78)
    {increasing sulfur content};

  % C — primary explanatory scene: localized trapping landscape.
  \node[panel letter] at (0.00,7.92) {C};
  \node[panel title] at (0.36,7.92) {Localized shallow and deep traps coexist};
  \begin{scope}[shift={(0.52,4.22)}]
    \draw[evidence axis] (0,0.00) -- (0,3.22);
    \node[small note, rotate=90] at (-0.34,1.62) {energy};
    \node[small note, anchor=west] at (0.10,2.91) {higher};
    \draw[cGray!42, line width=0.34pt, dashed]
      (0.42,2.48) -- (10.55,2.48);
    \node[small note, anchor=east, fill=white, inner sep=1.0pt]
      at (10.45,2.68) {transport level};

    \path[fill=cTeal!8, draw=none]
      (0.42,2.48)
      .. controls (1.20,2.36) and (1.58,2.18) .. (2.10,2.20)
      .. controls (2.52,2.22) and (2.62,1.58) .. (3.08,1.48)
      .. controls (3.54,1.58) and (3.66,2.25) .. (4.22,2.30)
      .. controls (5.02,2.35) and (5.50,2.02) .. (6.05,1.92)
      .. controls (6.46,1.78) and (6.47,0.58) .. (7.17,0.43)
      .. controls (7.88,0.58) and (7.82,1.94) .. (8.53,2.06)
      .. controls (9.18,2.14) and (9.62,2.32) .. (10.55,2.48)
      -- (10.55,0.02) -- (0.42,0.02) -- cycle;
    \draw[cTeal, line width=0.92pt]
      (0.42,2.48)
      .. controls (1.20,2.36) and (1.58,2.18) .. (2.10,2.20)
      .. controls (2.52,2.22) and (2.62,1.58) .. (3.08,1.48)
      .. controls (3.54,1.58) and (3.66,2.25) .. (4.22,2.30)
      .. controls (5.02,2.35) and (5.50,2.02) .. (6.05,1.92)
      .. controls (6.46,1.78) and (6.47,0.58) .. (7.17,0.43)
      .. controls (7.88,0.58) and (7.82,1.94) .. (8.53,2.06)
      .. controls (9.18,2.14) and (9.62,2.32) .. (10.55,2.48);

    \fill[cBlue] (3.08,1.72) circle (0.085);
    \fill[cBlue] (7.17,0.67) circle (0.085);
    \draw[cBlue!72, line width=0.42pt, dashed, -{Stealth[length=3.0pt,width=2.1pt]}]
      (3.08,1.83) -- (3.08,2.38);
    \draw[cBlue!72, line width=0.42pt, dashed, -{Stealth[length=3.0pt,width=2.1pt]}]
      (7.17,0.78) -- (7.17,2.38);
    \node[note, text=cTeal] at (3.08,1.05) {shallow trap};
    \node[note, text=cTeal] at (7.17,0.09) {deep trap};
    \node[small note, anchor=west] at (8.44,1.22)
      {lower state = deeper well};
  \end{scope}

  % Narrative links into the explanatory scene and evidence.
  \draw[causal] (2.05,8.62) to[out=-90,in=150] (2.76,7.68);
  \draw[causal] (6.62,8.52) to[out=-90,in=70] (6.45,7.56);

  % D — transient-current kinetic evidence.
  \node[panel letter] at (0.00,3.45) {D};
  \node[panel title] at (0.36,3.45) {Transient-current kinetics};
  \begin{scope}[shift={(0.58,0.42)}]
    \draw[evidence axis] (0,0) -- (4.55,0);
    \draw[evidence axis] (0,0) -- (0,2.42);
    \node[small note] at (4.30,-0.31) {time};
    \node[small note, rotate=90] at (-0.35,2.16) {$I(t)$};
    \draw[cBlue, line width=0.92pt]
      plot[smooth] coordinates {
        (0.18,2.15) (0.48,1.82) (0.88,1.48) (1.38,1.14)
        (2.05,0.82) (2.82,0.58) (3.65,0.41) (4.30,0.32)
      };
    \node[small note, anchor=west, text=cBlue] at (2.55,1.40)
      {trap-mediated decay};
    \node[small note, anchor=west] at (0.20,0.31) {constant applied $V$};
  \end{scope}

  % E — noncontact surface-potential-derived trap evidence.
  \node[panel letter] at (5.78,3.45) {E};
  \node[panel title] at (6.14,3.45) {ISPD-derived trap distribution};
  \begin{scope}[shift={(6.30,0.42)}]
    \draw[evidence axis] (0,0) -- (4.90,0);
    \draw[evidence axis] (0,0) -- (0,2.42);
    \node[small note] at (4.55,-0.31) {trap energy};
    \node[small note, rotate=90] at (-0.38,1.92) {$g(E)$};
    \path[fill=cRed!13, draw=none]
      (0.18,0.05)
      .. controls (0.72,0.12) and (1.03,0.62) .. (1.48,1.04)
      .. controls (1.93,1.52) and (2.34,1.86) .. (2.82,1.72)
      .. controls (3.34,1.55) and (3.68,0.86) .. (4.58,0.12)
      -- (4.58,0.05) -- cycle;
    \draw[cRed, line width=0.92pt]
      (0.18,0.05)
      .. controls (0.72,0.12) and (1.03,0.62) .. (1.48,1.04)
      .. controls (1.93,1.52) and (2.34,1.86) .. (2.82,1.72)
      .. controls (3.34,1.55) and (3.68,0.86) .. (4.58,0.12);
    \draw[cGray!45, line width=0.30pt, dashed] (2.82,0) -- (2.82,1.72);
    \node[small note, anchor=west] at (0.22,0.34) {shallower};
    \node[small note, anchor=east] at (4.52,0.34) {deeper};
    \node[small note, anchor=west, text=cRed] at (3.05,1.70)
      {derived evidence};
  \end{scope}

  \draw[causal] (4.12,4.02) to[out=-90,in=90] (3.10,3.62);
  \draw[causal] (7.50,4.02) to[out=-90,in=90] (8.45,3.62);

  % F — curated floating cantilever motif with caller-owned semantic overlays.
  \node[panel letter] at (12.34,10.83) {F};
  \node[panel title] at (12.70,10.83) {Floating cantilever response};
  \PanelFFloatingCantilever{f}{(15.34,9.55)}
  \node[small note, anchor=east, text=cAmber] at (14.83,8.98)
    {floating cantilever};
  \node[small note, anchor=west] at (16.98,8.17) {driven electrode};
  \node[small note, anchor=west, text=cBlue] at (17.08,10.25)
    {grounded source return};
  \draw[force arrow]
    ($(fTrappedCharge)+(-0.12,-0.42)$) -- ++(-1.05,-0.18);
  \node[small note, anchor=south, text=cRed]
    at ($(fTrappedCharge)+(-0.78,-0.46)$) {Coulomb force};
  \draw[causal]
    (11.22,5.78) to[out=10,in=190] (13.34,7.70);
  \node[small note, fill=white, inner sep=1.2pt]
    at (12.35,6.83) {trapped charge\to mechanical response};

  \node[small note, anchor=east, text=cGray!80]
    at (18.55,-0.15) {schematic concepts; D--E encode experimental evidence};
\end{tikzpicture}
\end{document}
